\chapter{Conclusão}
\label{chap:conclusao}

\section{Síntese das Contribuições}
\label{sec:sintese}

Esta dissertação apresentou o Fleet UI, um framework open-source de orquestração
de experimentos de repetibilidade de navegação autônoma para robôs móveis em
ROS~2, com arquitetura preparada para operação em frota. O sistema formaliza um
protocolo record--replay que transforma uma prática informal e fragmentada ---
executar um robô várias vezes e comparar os resultados ``na mão'' --- em um
pipeline documentado, automatizável e transferível.

As contribuições se organizam em três níveis:

\textbf{Nível arquitetural}: a separação de papéis MUUT/FUUT formaliza a distinção
entre unidades que executam tarefas e unidades que observam, permitindo designs
experimentais heterogêneos sem modificação do core do sistema. Os buffers TF por
robô, o executor multithreaded e os perfis de QoS diferenciados resolvem problemas
concretos de compatibilidade que qualquer equipe que adotar ROS~2 em modo multi-robô
enfrentará --- ainda que, como reconhecido na Seção~\ref{sec:limitacoes}, a
avaliação quantitativa desta dissertação não tenha exercitado esses mecanismos
sob operação simultânea de múltiplos robôs.

\textbf{Nível experimental}: o protocolo record--replay com análise RMSE pairwise
por reamostragem temporal normalizada constitui uma metodologia de avaliação que
pode ser adotada por outros grupos sem modificação. A campanha final em
simulação, composta por uma baseline limpa e $N=10$ replays independentes da
rota \texttt{dissertation\_\allowbreak clean01}, obteve RMSE médio de
3,35\,cm contra a execução de referência, RMSE máximo de 4,76\,cm, erro final
médio de 1,64\,cm e duração média de 17,51\,s. Esses resultados fornecem um
baseline empírico controlado para comparação com hardware real e com outras
pilhas de navegação. Campanhas exploratórias executadas antes do controle
rigoroso de pose inicial são reportadas apenas como achados diagnósticos, pois
demonstraram que não reinicializar o estado entre replays pode contaminar
diretamente a interpretação do RMSE.

\textbf{Nível de acessibilidade}: a interface web reduz a necessidade de
interação direta com a linha de comando do ROS~2 para configurar, executar e
monitorar experimentos. A visualização do mapa SLAM ao vivo com coordenadas
reais e o botão \textit{Executar record} representam uma redução perceptível
no esforço de configuração de um experimento do zero --- uma percepção
qualitativa dos autores, já que, como discutido na Seção~\ref{sec:limitacoes},
nenhum estudo de usabilidade controlado foi conduzido para quantificar essa
redução.

\section{Retomada das Perguntas de Pesquisa}
\label{sec:retomada_pp}

A Seção~\ref{sec:perguntas_pesquisa} formalizou quatro perguntas de pesquisa
(PP1--PP4), cada uma associada a uma hipótese (H1--H4) e avaliada por meio de
uma ou mais questões de avaliação (QA1--QA4, Seção~\ref{sec:questoes}). A
Tabela~\ref{tab:pp_h_qa_sintese} consolida essa correspondência e o resultado
de cada uma; o texto que segue detalha cada linha.

\begin{table}[H]
\centering
\caption{Correspondência entre perguntas de pesquisa, hipóteses, questões de avaliação e resultado.}
\label{tab:pp_h_qa_sintese}
\begin{tabularx}{\textwidth}{lllX}
\toprule
\textbf{PP} & \textbf{Hipótese} & \textbf{Avaliada por} & \textbf{Resultado} \\
\midrule
PP1 & H1 & QA1 & Corroborada \\
PP2 & H2 & QA2, QA3 & Corroborada (campanha final, $N=10$) \\
PP3 & H3 & QA4 & Evidência exploratória compatível, causalidade não estabelecida \\
PP4 & H4 & --- (qualitativa) & Evidência funcional apenas; esforço percebido não avaliado \\
\bottomrule
\end{tabularx}
\end{table}

\textbf{PP1 (automação de protocolo)}: corroborada. Os sete serviços do
orquestrador (Seção~\ref{sec:orquestrador}), o script
\texttt{experiment\_\allowbreak repeatability.py} (Seção~\ref{sec:impl_pipeline}) e o script de
campanha com relançamento do stack expressam a sessão de validação final sem
código adicional específico de experimento --- apenas parâmetros de linha de
comando (rota, número de réplicas, tópicos de coleta e diretório de saída).

\textbf{PP2 (repetibilidade mensurável)}: corroborada no cenário controlado
final. A campanha \texttt{dissertation\_\allowbreak clean01\_\allowbreak final\_\allowbreak manual},
com $N=10$ replays independentes iniciados do mesmo estado, apresentou RMSE
médio de 3,35\,cm e máximo de 4,76\,cm contra a execução de referência, além de
erro final médio de 1,64\,cm. Esses valores ficam abaixo do limiar prático de
25\,cm definido pela tolerância padrão de objetivo do Nav2.

\textbf{PP3 (persistência fora do regime controlado)}: nem corroborada nem
refutada de forma conclusiva. As campanhas exploratórias discutidas na
Seção~\ref{sec:res_exp_complementares} mostram que execuções sem controle de
pose inicial podem produzir métricas contaminadas por deslocamentos acumulados,
mas esses dados não permitem atribuir a variabilidade a uma causa única. PP3
permanece uma questão em aberto, a ser respondida por desenhos experimentais
dedicados com variação controlada de pose inicial, ação de navegação e estado do
SLAM.

\textbf{PP4 (custo de engenharia)}: há evidência funcional de redução da
interação manual, mas a redução de esforço percebido ou mensurável não foi
empiricamente avaliada. Nenhuma das execuções reportadas neste trabalho exigiu
escrever ou modificar código Python além dos parâmetros de linha de comando do
protocolo; a interface web descrita na Seção~\ref{sec:impl_web} permite o mesmo
fluxo sem terminal --- isso é verificável diretamente no código e nos fluxos
descritos no Apêndice~\ref{ap:cli_reference}. O que não foi medido é se essa
redução de interação manual se traduz em menor esforço percebido pelo
pesquisador: nenhum estudo de usabilidade com participantes externos (tempo de
configuração, número de erros, SUS, NASA-TLX ou instrumento equivalente) foi
conduzido --- uma lacuna explicitamente reconhecida na Seção~\ref{sec:limitacoes}
e retomada como trabalho futuro.

\section{Reprodutibilidade vs.\ Repetibilidade: Retomada}
\label{sec:repro_vs_repet}

A Seção~\ref{sec:motivacao} distinguiu repetibilidade (o mesmo sistema, no
mesmo ambiente, produzindo resultados equivalentes em execuções sucessivas) de
reprodutibilidade (um investigador independente, em outro ambiente, obtendo
resultados equivalentes a partir do artefato publicado). Essa distinção se
confirma, e não se dilui, à luz dos resultados apresentados: os dados da
campanha final controlada (Capítulo~\ref{chap:resultados}) medem repetibilidade
--- nenhum deles envolveu replicação por outro pesquisador ou laboratório. O
framework favorece reprodutibilidade por construção (código aberto, protocolo
documentado, artefatos estruturados de saída, Seção~\ref{sec:rel_reprodutibilidade}),
mas essa é uma propriedade de \emph{desenho}, não um resultado
\emph{demonstrado} nesta dissertação. Essa distinção não invalida a
contribuição do trabalho; delimita, com precisão, o que foi de fato
demonstrado empiricamente, em contraste com o que é apenas facilitado pelo
desenho do sistema --- e é por isso que a replicação independente do protocolo
por outro laboratório aparece explicitamente como primeiro passo de trabalho
futuro (Seção~\ref{sec:futuros}).

\section{Limitações}
\label{sec:limitacoes}

Além das limitações de validação em simulação e de ambiente estático discutidas
a seguir, as seguintes limitações adicionais merecem registro explícito por
afetarem diretamente a interpretação dos resultados desta dissertação:

\begin{itemize}
    \item \textbf{Avaliação quantitativa restrita a uma única unidade móvel.}
    Embora a arquitetura do Fleet UI suporte múltiplas unidades operando
    simultaneamente sob papéis distintos (MUUT/FUUT/SU, Seção~\ref{sec:papeis}),
    toda a avaliação quantitativa de repetibilidade reportada nos
    Capítulos~\ref{chap:avaliacao} e~\ref{chap:resultados} foi realizada com um
    único TurtleBot4 Standard. Escalabilidade e comportamento sob operação
    simultânea de múltiplos robôs --- incluindo eventuais interferências de
    rede, DDS ou TF entre unidades --- não foram empiricamente avaliados.

    \item \textbf{Escopo restrito da campanha quantitativa.} A campanha final
    utiliza $N=10$ replays, superando a limitação inicial de amostra muito
    pequena para o cenário controlado. No entanto, esse $N$ maior foi aplicado a
    uma única rota curta, em um único ambiente e com uma única plataforma
    simulada. Assim, a estimativa de variabilidade é mais robusta para o cenário
    \texttt{dissertation\_\allowbreak clean01}, mas não generaliza automaticamente
    para rotas longas, curvas acentuadas, obstáculos dinâmicos, múltiplos robôs
    ou hardware real.
\end{itemize}

A validação em simulação, embora controlada, não captura todos os fenômenos
relevantes em ambientes físicos: slip de roda em superfícies irregulares, variação
de temperatura e umidade que afeta a calibração do LiDAR, e deriva de IMU de longo
prazo. Experimentos em hardware real são a extensão mais urgente deste trabalho.

O protocolo atual assume que o ambiente permanece estático entre as execuções.
Qualquer variação --- um objeto movido, uma pessoa cruzando o campo do LiDAR ---
invalida a comparação entre execuções. A integração de detecção de mudanças ambientais
como pré-condição para aceitar uma execução de replay é um desenvolvimento natural.

\section{Trabalhos Futuros}
\label{sec:futuros}

A agenda de trabalhos futuros pode ser organizada em três horizontes.

\textbf{Curto prazo}: ampliar a campanha controlada para múltiplas rotas,
incluindo trajetórias mais longas, curvas fechadas e obstáculos dinâmicos;
repetir o protocolo em hardware físico com TurtleBot4 real; comparar
\texttt{/\allowbreak odom}, pose estimada por SLAM e \textit{ground truth} do
Gazebo em uma mesma campanha para separar erro de controle, erro de localização
e erro de medição; avaliar quantitativamente cenários multi-robô (MUUT+FUUT
simultâneos), respondendo diretamente à limitação de escopo single-robot
registrada na Seção~\ref{sec:limitacoes}; integrar câmera OAK-D para coleta de
dados visuais sincronizados com trajetória.

\textbf{Médio prazo}: extensão do protocolo para detectar e compensar variações
ambientais entre execuções; interface para exportação de dados no formato Roboflow
ou HuggingFace, facilitando o uso dos bags como dataset para treinamento de modelos
de navegação; suporte a experimentos com obstáculos dinâmicos programáveis; estudo
de usabilidade controlado com participantes externos para quantificar a redução de
esforço de engenharia atribuída à interface web (PP4).

\textbf{Longo prazo}: generalização do framework para robôs além da família
TurtleBot, incluindo manipuladores móveis; integração com plataformas de CI/CD
para robótica (\textit{continuous integration} de algoritmos de navegação via
execução automatizada do protocolo de repetibilidade); publicação do framework
no ecossistema ROS~2 como pacote instalável via \texttt{apt}; replicação
independente do protocolo por outro laboratório, como primeiro passo em direção
à reprodutibilidade no sentido forte discutido na Seção~\ref{sec:repro_vs_repet}.

\section{Consideração Final}
\label{sec:final}

A repetibilidade em robótica experimental não depende apenas de executar o mesmo
comando várias vezes; depende de controlar cuidadosamente o estado inicial, a
pilha de navegação, a coleta de dados e o método de comparação. A campanha final
desta dissertação mostra que, com infraestrutura adequada, é possível executar
dez replays independentes em simulação, todos partindo praticamente do mesmo
estado em \texttt{/\allowbreak odom}, e obter RMSE médio de 3,35\,cm em relação
à execução de referência. Mais importante que o número isolado é o protocolo que
o torna interpretável: baseline limpa, relançamento controlado entre replays,
coleta padronizada em MCAP e análise automatizada.
